% Physical pages 3 and 4 of the concise study: the movement of the formulary,
% stated as an argument and not as a summary of what follows. The developed
% comparison of the three readings begins on the fifth page.

\section{The Propers: Themes and Movement}
\zlabel{triptych:concise:themes:start}

This Mass asks God for life and is answered with his word. It opens with a
people in exile confessing that it has disobeyed God's commandments, sets at
its centre a father who asks Christ to come down to his dying son and is sent
home with a word instead, and ends by asking that those who have received the
sacred gifts may always obey the commandments the opening confessed broken.
Between confession and obedience stand a road, a song and a hope, and the hope
rests on a word that has been spoken and not yet seen fulfilled.

The Introit is cut from the prayer that Azarias prayed in the furnace at
Babylon (Dan 3:25--45). In the Missal's form it confesses first and asks
after: all that God has done to his people he has done \latin{in vero
iudicio}, in true judgment, \latin{quia peccavimus tibi, et mandatis tuis non
oboedivimus}; then, \latin{sed da gloriam nomini tuo}, glorify thy name and
deal with us according to the multitude of thy mercy. The prayer it is cut
from belongs to a people that has lost ``prince, or leader, or prophet, or
holocaust, or sacrifice'' (3:38) and offers instead ``a contrite heart and
humble spirit''. The priest says those words at every Mass when he has offered
the chalice, \latin{In spiritu humilitatis et in animo contrito suscipiamur a
te, Domine} (\work{Ordo Missae}, no.~1033). The verse after the antiphon is the
first of the longest psalm: ``Blessed are the undefiled in the way, who walk in
the law of the Lord.'' A people that confesses it has not obeyed is shown at
once the road of obedience.

The Collect asks two gifts of a God \latin{placatus}, as one appeased:
\latin{indulgentiam} and \latin{pacem}, pardon and then peace, so that the
faithful may be cleansed from every offence and serve him \latin{secura mente},
with a mind at rest. The Epistle begins with the verse's own verb,
\latin{Videte quomodo caute ambuletis}, walk carefully, \latin{non quasi
insipientes, sed ut sapientes}. St Paul writes it to Gentiles who were
``heretofore darkness, but now light in the Lord'' (Eph 5:8), and tells them
how to spend their time in a hostile world: ``redeeming the time, because the
days are evil''; understanding ``what is the will of God''; filled not with
wine but with the Holy Spirit; \latin{cantantes, et psallentes in cordibus
vestris Domino}, singing and making melody in their hearts to the Lord; and
``Giving thanks always for all things''. Its last verse, ``Being subject
one to another, in the fear of Christ'', turns the letter to the household.

Between the lessons the Gradual lifts \latin{Oculi omnium}, the eyes of all,
to the God who gives them
food \latin{in tempore opportuno}, in due season, and fills \latin{omne
animal}, every living creature, with blessing; the psalm goes on to say that
``The Lord is just in all his ways'' (Ps 144:17), the confession with which the
Introit began. The Alleluia sings the heart the Epistle asked for, and in its
two verbs: \latin{Paratum cor meum, Deus, paratum cor meum: cantabo, et psallam
tibi, gloria mea}. Where the Vulgate reads \latin{psallam in gloria mea}, the
Missal addresses the song to God and names him the singer's glory.

The Gospel is the second sign that Jesus did in Galilee. St John sets it after
two days among the Samaritans, of whom ``many more believed in him, because of
his own word'' (Jn 4:41), and after the saying ``that a prophet hath no honour
in his own country'' (4:44). A royal official, \latin{regulus}, whose son
\latin{incipiebat mori}, was beginning to die at Capharnaum, comes to Jesus and
asks him to come down and heal him. Jesus answers with a rebuke:
\latin{Nisi signa et prodigia videritis, non creditis}. The
father asks again, \latin{Domine, descende priusquam moriatur filius meus}, and
is given not a presence but a word: \latin{Vade, filius tuus vivit}.
\latin{Credidit homo sermoni, quem dixit ei Iesus, et ibat}: the man believed
the word and went. On the road his servants meet him; he asks the hour, and
they answer \latin{heri hora septima}, yesterday at the seventh hour; it was
the hour of the word, \latin{et credidit ipse, et domus eius tota}. The Missal
begins the lesson after the mention of Cana and ends it before the evangelist
counts the sign, so that it reads the healing alone. It moves from asking, through a word believed in absence, to a faith that reaches
a whole household.

The Offertory turns from the healed house back to exile. \latin{Super flumina
Babylonis illic sedimus, et flevimus: dum recordaremur tui, Sion}: by the rivers
of Babylon we sat and wept, remembering thee, Sion. The Missal's \latin{tui}
speaks to a city remembered and not seen. The psalm's next verses ask, ``How shall we sing the song of the Lord in a
strange land?'' (Ps 136:4), the question the Epistle and the Alleluia have
already answered, and end in a curse on Babylon's little ones that the Mass
does not sing (136:9). The Secret asks that \latin{haec mysteria} furnish
\latin{caelestem \dots\ medicinam}, heavenly medicine, and purge the vices of
the heart; its vocabulary meets the Gospel's sick son and the fever that left
him, and the mysteries are the subject of both its verbs.

The Communion returns to Ps~118, whose first verse the Introit sang:
\latin{Memento verbi tui servo tuo, Domine, in quo mihi spem dedisti: haec me
consolata est in humilitate mea}. Remember thy word, in which thou hast given
me hope; this has comforted me in my humiliation. Azarias had prayed
\latin{sumusque humiles in universa terra hodie propter peccata nostra}, we are
brought low in all the earth this day for our sins (Dan 3:37); the Communion
finds comfort in that same lowliness. The psalm gives the reason the Mass leaves
unsung, \latin{quia eloquium tuum vivificavit me}, because thy word has given
me life, and the Gospel's word was \latin{Filius tuus vivit}. The Postcommunion
asks that those who have received may be made worthy of the sacred gifts,
\latin{fac nos \dots\ tuis semper oboedire mandatis}. Its noun and its verb are
the Introit's, so that the Mass ends by asking for the obedience it began by
confessing it had not given.

Three threads run through the whole. The first is a word given in absence: the
promise the Communion recalls, the word Christ speaks instead of coming down,
the Sion that is remembered and not seen, and the mysteries in which the Secret
asks a cure. The second is the condition of those who pray: sinners and
exiles, in evil days and humiliation, fed in due season and not before. The
third is the voice they pray with: a confession and a lament that are first
Israel's, from the furnace and the rivers of Babylon, beside chants of all eyes
and every living creature and an Epistle to Gentiles made light.

Three readings of the whole Mass grow from three questions, each developed at
length in the full study with four senses of its own. The first asks what faith
Christ asks of the father and why he heals by a word without going down. St
Gregory the Great answers that the father sought the bodily presence of one who
was absent nowhere; St John Chrysostom, that Christ healed the father's mind
together with the son's body, teaching him to attend to his word and not to his
signs. The reading hears the Introit's confession, the Collect's mind at rest,
the Gradual's waiting and the Communion's hope as the school of that faith. The
second asks how a people that has sinned and lives in exile is to pray, sing
and spend its time. It is heard from the Introit, the Offertory and the
Epistle: St Augustine reads the rivers of Babylon as all that is loved here and
passes, beside which the citizens of Jerusalem sit and weep; St Hilary of
Poitiers, as the soul's exile since Adam from the heavenly Jerusalem, its
mother. The third asks whose story the healing at Capharnaum tells and whose
voices the Mass lets us hear. St Augustine reads the episode in its chapter, as
a figure of the people of the Jews, who saw signs and scarcely believed, and of
the nations, who believe having seen none; St Thomas Aquinas reads the ruler as
one of the Fathers of the Old Testament praying for a sick people.

The three hold less in common than their shared Gospel suggests. All three
hear the Postcommunion's \latin{oboedire mandatis} answer the Introit's
\latin{non oboedivimus}, begin from the Introit's confession though each hears
it differently, and give the Communion's hope in God's word a place. The first
and third take the Gospel's word as the hinge of the Mass; the second gives the
Gospel a supporting place and centres on the chants and the Epistle. None of
them denies what another holds, and the differences inside them are real: how
much the ruler believed when he came, what the willows of Babylon are, whom the
ruler figures, and whose exile the Offertory first voices. The commentary that
follows sets the readings' answers beside one another at the points where the
texts put the same question to all three.
\zlabel{triptych:concise:themes:end}
